%-------------------------------------------------------------------------------
%	SECTION TITLE
%-------------------------------------------------------------------------------
\cvsection{项目\! 经历}


%-------------------------------------------------------------------------------
%	CONTENT
%-------------------------------------------------------------------------------
\begin{cventries}
%---------------------------------------------------------
\cventry
  {软件工程师（项目合作）}
  {苏州新苏报业网络科技有限公司}
  {中国，江苏省，苏州市}
  {2026 年 4 月 -- 2026 年 8 月}
  {
    \begin{cvitems}
      \item {\textbf{新闻校对智能体（4--5 月，独立开发）：}完成需求拆解、架构设计及端到端交付；构建 Vue 编辑器 SDK 与 FastAPI/SQLAlchemy/LangChain 后端，实现 Agent 编排、事实检索及修改建议生成。}
      \item {本地部署 Qwen，设计异步任务与错误处理流程，并使用 Langfuse 跟踪模型调用、任务状态及端到端链路。}
      \item {\textbf{工业图纸/电路图识别原型：}参与论文方法复现、数据标注及模型训练评估，覆盖 YOLO、UNet、HRNet 等检测与分割模型。}
      \item {\textbf{苏州新闻网管理后台（6--8 月，两人协作）：}负责基于 RuoYi-Vue-Plus（Spring Boot/Vue 3）的管理端前后端开发及门户端联调。}
      \item {实现多租户内容管理、视频压缩、旧系统数据迁移、第三方系统数据同步及静态站点定时构建调度。}
    \end{cvitems}
  }

%---------------------------------------------------------
\cventry
  {软件工程师（个人开源项目）}
  {\href{https://github.com/FreeFlyingSheep/stock-analysis}{A 股基本面分析 Agent}}
  {中国，江苏省，苏州市}
  {2026 年 2 月 -- 2026 年 3 月}
  {
    \begin{cvitems}
      \item {构建 A 股基本面评分与问答系统，基于 FastAPI、SQLAlchemy、PgQueuer 实现 CNInfo/Yahoo Finance 多源采集、异步调度及声明式 YAML 评分规则。}
      \item {构建 LangGraph + RAG + SSE 流式分析链路，集成 MCP、PostgreSQL/pgvector 与 MinIO；使用 DeepEval 建立 RAG、Agent 及端到端离线回归评估。}
      \item {通过 pytest/Testcontainers 验证服务依赖，提供 Docker Compose 与 Kustomize/Minikube 部署配置，并接入 OTel、Prometheus、Grafana、Loki、Tempo 可观测性栈。}
    \end{cvitems}
  }

%---------------------------------------------------------
\cventry
  {内核工程师}
  {系统软件与开源贡献}
  {中国，江苏省，苏州市}
  {2020 年 6 月 -- 2023 年 8 月}
  {
    \begin{cvitems}
      \item {\href{https://github.com/FreeFlyingSheep/valgrind-loongarch64}{为 Valgrind 添加 LoongArch64/Linux 支持}，完成核心组件移植、关键指令翻译与运行时验证，1 个月内达到初步可用并通过上百项回归测试。}
      \item {适配 Linux 内核、klibc、Kcov、RT-Thread 等组件并支持 LS2K0300 集成，覆盖底层库、测试工具、RTOS 与硬件平台。}
      \item {完成 OpenStack Nova/DevStack 的 Fedora MIPS 适配，1 个月内跑通核心服务链路。}
      \item {设计文档协作 CI/CD，支持实时 PR 预览与贡献者自动归因，提升审核及社区协作效率。}
    \end{cvitems}
  }

\end{cventries}
